% Generated by experiments/render_study3.py from experiments/results-v3; do not edit by hand.

\subsection{Study 3: generated mutants}
\label{sec:results-mutants}

Cosmic Ray generated 119 mutants of the 82-line subject. 55 raise an exception on every input, so no check is needed to see them. Of the 64 that run, 20 return the clean forecast on all 24 screening inputs and are counted as equivalent; 44 change a forecast. Of the 16 catalogue mutants one is equivalent, a target scaler reused across histories (the Ridge fit is linear in its targets, so the inverse transform undoes it), and none of the four Study 2 bugs is. The equivalent sources were also run on every evaluation case, and none differed from the clean pipeline there; 6 and 48 screening inputs give the same equivalent mutants as 24.

Table~\ref{tab:mutants} counts, for each check, the mutants caught in at least one of the six conditions and, in parentheses, the pairs of a mutant and a condition in which it was caught.

\begin{table}[htbp]\centering\small

\caption{Study 3, Ridge pipeline: mutants that change a forecast and are caught by each check in at least one of six conditions; in parentheses, the mutant--condition pairs caught. A deterministic check catches a mutant in a condition when it alarms on any of the 93 evaluation cases, a calibrated check when its Holm-adjusted test rejects (Section~\ref{sec:third-study}). The counts describe these sets of mutants, not a population of real faults.}

\label{tab:mutants}

\begin{tabular}{@{}lrrr@{}}\toprule

Check & Cosmic Ray & Catalogue & Study 2 bugs \\

 & (44 mutants, 264 pairs) & (15, 90) & (4, 24) \\ \midrule

Output checks & 3 (14) & 0 (0) & 0 (0) \\

Common shift check & 22 (132) & 2 (12) & 1 (6) \\

Per-series shift check & 22 (131) & 4 (24) & 2 (12) \\

Second implementation & 44 (256) & 15 (90) & 4 (24) \\

Calibrated loss alarm & 37 (186) & 13 (39) & 4 (16) \\

Answer-key alarm & 40 (210) & 14 (55) & 4 (20) \\

Paired check with the clean pipeline & 43 (215) & 15 (63) & 4 (22) \\

\bottomrule\end{tabular}\end{table}

The second implementation flagged every mutant in every condition in which the mutant returned a forecast. The 8 pairs it did not flag belong to 2 generated mutants that raise an exception on every case of four conditions, where the output checks report them. The paired check missed one generated mutant, which returns no valid forecast on any case and so leaves nothing to pair; the output checks caught it. The answer-key alarm caught at least as many mutants as the loss alarm in every set: 40 against 37 generated mutants (210 against 186 pairs), 14 against 13 catalogue mutants (55 against 39), and the Study 2 bugs in 20 against 16 pairs.

Each shift check caught 22 of the 44 generated mutants. Grouped by the function they change, a breakdown made after the results were known, the shift checks caught all 11 mutants of the forecasting loop, 11 of the 21 mutants of the two functions that apply and invert the scaling, and none of the 12 mutants of the lag construction and of the scaler fit. The catalogue agrees: the common shift check caught only its two inverse-transform mutants, the per-series check those and the two column swaps, and the other 11 (lag window, lag order, recursion state, horizon index, seasonal period, dropped training rows and a reused feature scaler) passed both.

No deterministic check alarmed on the clean pipeline, on 332 cases in each condition (239 calibration and 93 evaluation cases). At the 5\% level the loss alarm fired on 29/558 [3.6, 7.4]\% of the clean evaluation cases and the answer-key alarm on 33/558 [4.2, 8.2]\%.

\subsection{Study 3: learners beyond Ridge}
\label{sec:results-learners}

Table~\ref{tab:shift-classes} classifies each learner on its clean pipeline, on 16 cases in each of the six conditions. Every learner received the class that the protocol expected.

\begin{table}[htbp]\centering\small

\caption{Study 3: how each learner's clean pipeline responds to a shift of 10 in its history. The public shift check asks for the same forecast plus 10; the calibrated check asks only that the shifted forecast look like one more refit with a new training seed. A learner is exact when the public check passes on every case, in distribution when the calibrated check fails no more often than its 5\% level allows, and not applicable otherwise.}

\label{tab:shift-classes}

\begin{tabular}{@{}lrrl@{}}\toprule

Learner & Public: PASS & Calibrated: FAIL & Class \\ \midrule

Ridge with global scaling (control) & 96/96 & --- & exact \\

Histogram gradient boosting & 96/96 & --- & exact \\

Perceptron, window-standardized lags & 0/96 & 4/96 & in distribution \\

Perceptron, raw lags & 0/96 & 51/96 & not applicable \\

\bottomrule\end{tabular}\end{table}

The three learners run the subject without its scaling step, so a mutant of a scaling function changes nothing. Of the first 84 Cosmic Ray mutants in the seeded order, 24 raise an exception on every input and 48 are equivalent, 34 of them in the three scaling functions; the 12 that remain are the same for the three learners. Of the 16 catalogue mutants 4 are equivalent, and one of the four Study 2 bugs, as the protocol expected. Table~\ref{tab:learners} counts the mutants caught.

\begin{table}[htbp]\centering\small

\caption{Study 3: mutants caught in at least one condition with each learner, as Cosmic Ray / catalogue / Study 2 bugs (of 12 / 12 / 3). A dash marks a check that does not apply to the learner's shift class (Table~\ref{tab:shift-classes}); no second implementation exists for these learners.}

\label{tab:learners}

\begin{tabular}{@{}lccc@{}}\toprule

Check & Boosting & Perceptron, standardized & Perceptron, raw \\ \midrule

Output checks & 3 / 0 / 0 & 3 / 0 / 0 & 3 / 0 / 0 \\

Common shift check & 6 / 0 / 0 & --- & --- \\

Per-series shift check & 6 / 2 / 1 & --- & --- \\

Calibrated shift check & --- & 0 / 0 / 0 & --- \\

Calibrated loss alarm & 8 / 11 / 3 & 5 / 8 / 3 & 8 / 6 / 3 \\

Answer-key alarm & 9 / 12 / 3 & 5 / 7 / 3 & 8 / 7 / 3 \\

Paired check with the clean pipeline & 11 / 12 / 3 & 6 / 9 / 3 & 11 / 12 / 3 \\

\bottomrule\end{tabular}\end{table}

The calibrated shift check kept its level and caught no mutant. On the clean window-standardized pipeline it returned FAIL on 27/558 [3.3, 6.9]\% of the evaluation cases. On the 25 evaluated mutants it returned FAIL on 608 of 12,648 cases (4.8\%); 5 of the 136 pairs of a mutant and a condition have an unadjusted $p$-value of at most 0.05, no more than its level gives by chance, and none is left after Holm's correction. With this learner the paired check also caught fewer mutants than with the other two: 6 generated and 9 catalogue mutants, against 11 and 12 with each of the others. Every fit of this learner stopped at its limit of 500 iterations (22,878 convergence warnings on the 558 clean evaluation cases, 41 fits a case); the perceptron on raw lags did so on 888 of its 1,992 clean calibration and evaluation cases, and gradient boosting raised no warning. With gradient boosting the exact shift checks behaved as with Ridge: no alarm on its 1,992 clean cases, 6 of the 12 generated mutants caught, and at most 2 catalogue mutants.

No deterministic check alarmed on a clean pipeline (332 cases for each learner and condition). At the 5\% level the loss alarm fired on 30, 28 and 29 of the 558 clean evaluation cases of gradient boosting, the window-standardized perceptron and the raw-lag perceptron, and the answer-key alarm on 17, 32 and 22.

\subsection{Study 3: two questions about a suite}
\label{sec:results-suites}

Table~\ref{tab:suite-questions} repeats the covariance suites of Study 2 on 93 new suites with 239 calibration draws. Of the clean suites, 17 contain at least one unadjusted rejection, 2 a Holm rejection, and Fisher's combination flags 5, a rate whose interval holds the nominal 5\%. Under the partial shuffle all three flag every suite; on the condition with independent targets, where the shuffle changes nothing, the unadjusted check rejected in 4 suites and Holm in 3. Under the reduced spread Fisher's combination flags all 93 suites, while Holm rejects at least one condition in 68 and a given condition in 14 to 45.

\begin{table}[htbp]\centering\small

\caption{Study 3: suites of four covariance checks flagged, of 93, with Wilson 95\% intervals (percent). Raw counts a suite with any unadjusted rejection, Holm one with any rejection after Holm's adjustment, and Fisher one whose combined test rejects. Every flag on the control is a false alarm; the other two rows measure detection.}

\label{tab:suite-questions}

\begin{tabular}{@{}lrrr@{}}\toprule

Variant & Raw & Holm & Fisher \\ \midrule

Exact-reference control & 17/93 [11.7, 27.3] & 2/93 [0.6, 7.5] & 5/93 [2.3, 12.0] \\

Partial dependence shuffle & 93/93 [96.0, 100.0] & 93/93 [96.0, 100.0] & 93/93 [96.0, 100.0] \\

Reduced spread & 93/93 [96.0, 100.0] & 68/93 [63.3, 81.1] & 93/93 [96.0, 100.0] \\

\bottomrule\end{tabular}\end{table}

\subsection{Study 3: a joint-spread decision with a margin}
\label{sec:results-margin}

Table~\ref{tab:margin} gives how often each one-sided verdict was returned when the true distance lay exactly on the margin, where each was designed to occur in at most about 5\% of repetitions. FAIL meets the design: between 0 and 29 of 381 repetitions over the 32 cells, with a 95\% interval above 5\% in 1 of them. WITHIN\_TOLERANCE does not: between 5 and 111 of 381, with the interval above 5\% in 16 of the 32 cells, 11 of them with a spread that is too wide and 13 in the two conditions with dependence over time. With the distance 10\% beyond the margin it was still returned in up to 95 of 381 repetitions. The expectation that the protocol wrote down for this verdict was therefore wrong.

\begin{table}[htbp]\centering\footnotesize

\caption{Study 3: verdicts of the margin decision when the true distance equals the margin, as FAIL / WITHIN\_TOLERANCE counts of 381 repetitions; the remainder is INCONCLUSIVE. Either verdict is an error here, designed to occur in at most about 5\% of repetitions (19 of 381).}

\label{tab:margin}

\begin{tabular}{@{}llrrrr@{}}\toprule

 & & \multicolumn{2}{c}{Margin 0.25} & \multicolumn{2}{c}{Margin 0.4375} \\

Condition & Spread & 64 paths & 256 paths & 64 paths & 256 paths \\ \midrule

White noise: independent targets & narrower & 20 / 5 & 13 / 15 & 20 / 14 & 17 / 13 \\

 & wider & 1 / 20 & 3 / 26 & 0 / 21 & 6 / 32 \\

\addlinespace

White noise: correlated targets & narrower & 12 / 9 & 13 / 18 & 20 / 8 & 24 / 20 \\

 & wider & 1 / 26 & 4 / 32 & 1 / 28 & 12 / 23 \\

\addlinespace

Negative AR(1) & narrower & 11 / 30 & 13 / 30 & 25 / 20 & 16 / 18 \\

 & wider & 0 / 53 & 0 / 52 & 1 / 62 & 5 / 47 \\

\addlinespace

Correlated random walk with drift & narrower & 10 / 43 & 17 / 43 & 29 / 31 & 20 / 12 \\

 & wider & 0 / 104 & 2 / 78 & 1 / 111 & 4 / 58 \\

\bottomrule\end{tabular}\end{table}

The estimator itself shows no bias: over the 240 cells its mean differs from the squared distance by at most 2.7 of its standard errors, and by more than two in 7. A reading made after the run locates the problem in the interval. The jackknife standard error has the right size on average (its mean is 1.05 times the standard deviation of the estimator in the median cell), but with a spread that is too wide it moves with the estimate (correlation between 0.85 and 0.98 in the 16 such cells of Table~\ref{tab:margin}), so a low estimate arrives with a short interval and falls below the margin. At a distance of zero, where it is the correct verdict, WITHIN\_TOLERANCE was returned in 670 of 1,524 repetitions at 64 paths and 1,482 at 256 with the margin 0.25, and in 1,443 and 1,524 with the margin 0.4375.

\subsection{Study 3: predictability estimates against the exact floor}
\label{sec:results-predictability}

Table~\ref{tab:predictability} divides the spectral bound by the exact unavoidable error on 200 series of each scorecard cell. On the eight stationary rows (pure noise, a fixed level, mean reversion and irregular cycles) the median ratio lies between 0.86 and 0.98, as the protocol expected. On the four rows with a trend it is far above 1 (24.00 and 24.28 at low noise, 2.35 and 2.47 at high noise), as expected. Two expectations were wrong: on the seasonal rows (1.12 and 1.03) and on the level-shift rows (0.83 and 0.83) the bound is close to the exact value, where the protocol expected it to be off. Those ratios use the realized future segment, as the bound is defined; computed from the history alone, the level-shift ratio is 81.1 at low noise. The random-walk rows have no stationary spectrum, so the bound is not defined there in theory; its values are reported as computed. Weighted permutation entropy is highest on the pure-noise and fixed-level rows. Across the 20 rows its rank correlation with the exact floor is 0.26, and 0.72 once the floor is divided by the variance of the history, which removes the scale that the entropy ignores.

\begin{table}[htbp]\centering\small

\caption{Study 3: the spectral bound divided by the exact unavoidable error (median and 5th to 95th percentile over 200 series; 1 is exact), and the median weighted permutation entropy (0 to 1), at low and high noise. The two random-walk rows have no stationary spectrum.}

\label{tab:predictability}

\begin{tabular}{@{}lrrrr@{}}\toprule

 & \multicolumn{2}{c}{Spectral bound / exact floor} & \multicolumn{2}{c}{Entropy} \\

\cmidrule(lr){2-3}\cmidrule(lr){4-5}

Kind of series & Low noise & High noise & Low & High \\ \midrule

Pure noise & 0.86 [0.52, 1.47] & 0.86 [0.52, 1.47] & 0.99 & 0.99 \\

Noise around a fixed level & 0.86 [0.52, 1.47] & 0.86 [0.52, 1.47] & 0.99 & 0.99 \\

Straight-line trend & 24.00 [19.88, 29.77] & 2.35 [1.20, 4.02] & 0.95 & 0.99 \\

Repeating seasonal cycle & 1.12 [0.65, 1.79] & 1.03 [0.60, 1.67] & 0.63 & 0.97 \\

Trend plus seasonal cycle & 24.28 [20.00, 28.10] & 2.47 [1.42, 3.60] & 0.63 & 0.97 \\

Sudden level shift & 0.83 [0.50, 1.34] & 0.83 [0.50, 1.34] & 0.87 & 0.99 \\

Mean reversion (AR(1)) & 0.98 [0.32, 3.64] & 0.94 [0.43, 1.90] & 0.94 & 0.99 \\

Irregular cycles (AR(2)) & 0.98 [0.21, 2.89] & 0.92 [0.39, 2.35] & 0.64 & 0.97 \\

Random walk & 0.85 [0.15, 3.88] & 0.83 [0.33, 2.38] & 0.97 & 0.99 \\

Random walk with drift & 2.46 [0.20, 8.86] & 1.52 [0.38, 4.68] & 0.95 & 0.99 \\

\bottomrule\end{tabular}\end{table}

\subsection{Study 3: scorecards of nine models}
\label{sec:results-zoo}

Table~\ref{tab:zoo} gives the scorecard of each model beside the row's attainable mark. No check failed for any model. The scores show where each model stops; they rank no model, because the rows are mechanisms chosen for their known answers and not a sample of real series. At low noise, ETS scores 35 on irregular cycles (mark 93); gradient boosting on lags scores 5 on a straight-line trend (mark 95), which a tree cannot extrapolate, and 93 on a sudden level shift (mark 97), where Theta scores 22; the perceptron on lags scores 4 on a random walk (mark 100).

Chronos is the one pretrained model. Before its run the protocol declared how each kind of series relates to the synthetic part of its pretraining data, which is drawn from Gaussian processes with constant, white-noise, linear, RBF, rational-quadratic and periodic kernels~\cite{ansari2024chronos}: five kinds resemble it (pure noise, a fixed level, a straight-line trend, a seasonal cycle and a trend with a seasonal cycle), irregular cycles resemble it in part, and four are not in its kernel bank (a level shift, mean reversion and the two random walks). Chronos scores 23 to 99 on the five kinds that resemble that data, 73 and 78 on irregular cycles, and 59 to 97 on the four that it does not contain. Its lowest score, 23, is on the low-noise straight-line trend (mark 95), and it scores 92 and 97 on mean reversion (marks 93 and 96). The declared resemblance does not order its scores; what the model saw in its real pretraining data cannot be stated.

\begin{table}[htbp]\centering\footnotesize\setlength{\tabcolsep}{3.4pt}

\caption{Study 3: scores out of 100 of nine models on the default scorecard grid (30 series a cell) at low and high noise, beside the attainable mark. Last is the last value, Seas.\ the seasonal naive forecast (period 12), Boost.\ and Perc.\ gradient boosting and a perceptron on lags, Chronos the pretrained small model, Parrot context parroting. The interval of every score is in the study's records. No ranking is claimed.}

\label{tab:zoo}

\begin{tabular}{@{}p{0.24\linewidth}lrrrrrrrrrr@{}}\toprule

Kind of series & Noise & Mark & Last & Seas. & ETS & ARIMA & Theta & Boost. & Perc. & Chronos & Parrot \\ \midrule

Pure noise & low & 99 & 48 & 47 & 99 & 99 & 95 & 79 & 37 & 99 & 47 \\

 & high & 99 & 48 & 47 & 99 & 99 & 95 & 79 & 54 & 99 & 47 \\

Noise around a fixed level & low & 99 & 48 & 47 & 99 & 99 & 96 & 79 & 66 & 97 & 47 \\

 & high & 99 & 48 & 47 & 99 & 99 & 96 & 79 & 23 & 95 & 47 \\

Straight-line trend & low & 95 & 8 & 6 & 96 & 92 & 76 & 5 & 21 & 23 & 6 \\

 & high & 95 & 31 & 33 & 94 & 96 & 87 & 44 & 57 & 48 & 29 \\

Repeating seasonal cycle & low & 98 & 3 & 47 & 90 & 84 & 87 & 65 & 47 & 73 & 49 \\

 & high & 98 & 22 & 47 & 90 & 78 & 87 & 71 & 47 & 78 & 42 \\

Trend plus seasonal cycle & low & 94 & 2 & 6 & 86 & 81 & 87 & 4 & 45 & 35 & 5 \\

 & high & 94 & 25 & 36 & 86 & 72 & 86 & 35 & 50 & 63 & 26 \\

Sudden level shift & low & 97 & 49 & 49 & 70 & 69 & 22 & 93 & 6 & 85 & 47 \\

 & high & 97 & 49 & 49 & 90 & 86 & 69 & 85 & 44 & 83 & 47 \\

Mean reversion (AR(1)) & low & 93 & 62 & 51 & 65 & 91 & 63 & 70 & 52 & 92 & 50 \\

 & high & 96 & 53 & 51 & 84 & 92 & 81 & 73 & 62 & 97 & 51 \\

Irregular cycles (AR(2)) & low & 93 & 35 & 30 & 35 & 94 & 35 & 55 & 51 & 73 & 54 \\

 & high & 93 & 43 & 37 & 45 & 88 & 45 & 66 & 54 & 78 & 50 \\

Random walk & low & 100 & 97 & 75 & 97 & 92 & 90 & 88 & 4 & 89 & 58 \\

 & high & 99 & 70 & 64 & 91 & 88 & 91 & 79 & 40 & 87 & 56 \\

Random walk with drift & low & 87 & 60 & 44 & 73 & 83 & 85 & 43 & 27 & 59 & 41 \\

 & high & 90 & 57 & 50 & 78 & 85 & 87 & 58 & 62 & 67 & 42 \\

\bottomrule\end{tabular}\end{table}
